\documentclass{article}

\usepackage{arxiv}

\usepackage[utf8]{inputenc} % allow utf-8 input
\usepackage[T1]{fontenc}    % use 8-bit T1 fonts
\usepackage{lmodern}        % https://github.com/rstudio/rticles/issues/343
\usepackage{hyperref}       % hyperlinks
\usepackage{url}            % simple URL typesetting
\usepackage{booktabs}       % professional-quality tables
\usepackage{amsfonts}       % blackboard math symbols
\usepackage{nicefrac}       % compact symbols for 1/2, etc.
\usepackage{microtype}      % microtypography
\usepackage{graphicx}

\title{Prediction of Adverse Drug Reactions - Using Advanced Statistical
Methods}

\author{
    Arshad Ahamed Shajahan
   \\
    School of Mathematics \\
    University of Galway \\
  Galway, Ireland \\
  \texttt{\href{mailto:A.Shajahan2@universityofgalway.ie}{\nolinkurl{A.Shajahan2@universityofgalway.ie}}} \\
   \And
    Sai Shasank Dandu
   \\
    School of Mathematics \\
    University of Galway \\
  Galway, Ireland \\
  \texttt{\href{mailto:S.Dandu2@universityofgalway.ie}{\nolinkurl{S.Dandu2@universityofgalway.ie}}} \\
  }

% Pandoc syntax highlighting
\usepackage{color}
\usepackage{fancyvrb}
\newcommand{\VerbBar}{|}
\newcommand{\VERB}{\Verb[commandchars=\\\{\}]}
\DefineVerbatimEnvironment{Highlighting}{Verbatim}{commandchars=\\\{\}}
% Add ',fontsize=\small' for more characters per line
\usepackage{framed}
\definecolor{shadecolor}{RGB}{248,248,248}
\newenvironment{Shaded}{\begin{snugshade}}{\end{snugshade}}
\newcommand{\AlertTok}[1]{\textcolor[rgb]{0.94,0.16,0.16}{#1}}
\newcommand{\AnnotationTok}[1]{\textcolor[rgb]{0.56,0.35,0.01}{\textbf{\textit{#1}}}}
\newcommand{\AttributeTok}[1]{\textcolor[rgb]{0.13,0.29,0.53}{#1}}
\newcommand{\BaseNTok}[1]{\textcolor[rgb]{0.00,0.00,0.81}{#1}}
\newcommand{\BuiltInTok}[1]{#1}
\newcommand{\CharTok}[1]{\textcolor[rgb]{0.31,0.60,0.02}{#1}}
\newcommand{\CommentTok}[1]{\textcolor[rgb]{0.56,0.35,0.01}{\textit{#1}}}
\newcommand{\CommentVarTok}[1]{\textcolor[rgb]{0.56,0.35,0.01}{\textbf{\textit{#1}}}}
\newcommand{\ConstantTok}[1]{\textcolor[rgb]{0.56,0.35,0.01}{#1}}
\newcommand{\ControlFlowTok}[1]{\textcolor[rgb]{0.13,0.29,0.53}{\textbf{#1}}}
\newcommand{\DataTypeTok}[1]{\textcolor[rgb]{0.13,0.29,0.53}{#1}}
\newcommand{\DecValTok}[1]{\textcolor[rgb]{0.00,0.00,0.81}{#1}}
\newcommand{\DocumentationTok}[1]{\textcolor[rgb]{0.56,0.35,0.01}{\textbf{\textit{#1}}}}
\newcommand{\ErrorTok}[1]{\textcolor[rgb]{0.64,0.00,0.00}{\textbf{#1}}}
\newcommand{\ExtensionTok}[1]{#1}
\newcommand{\FloatTok}[1]{\textcolor[rgb]{0.00,0.00,0.81}{#1}}
\newcommand{\FunctionTok}[1]{\textcolor[rgb]{0.13,0.29,0.53}{\textbf{#1}}}
\newcommand{\ImportTok}[1]{#1}
\newcommand{\InformationTok}[1]{\textcolor[rgb]{0.56,0.35,0.01}{\textbf{\textit{#1}}}}
\newcommand{\KeywordTok}[1]{\textcolor[rgb]{0.13,0.29,0.53}{\textbf{#1}}}
\newcommand{\NormalTok}[1]{#1}
\newcommand{\OperatorTok}[1]{\textcolor[rgb]{0.81,0.36,0.00}{\textbf{#1}}}
\newcommand{\OtherTok}[1]{\textcolor[rgb]{0.56,0.35,0.01}{#1}}
\newcommand{\PreprocessorTok}[1]{\textcolor[rgb]{0.56,0.35,0.01}{\textit{#1}}}
\newcommand{\RegionMarkerTok}[1]{#1}
\newcommand{\SpecialCharTok}[1]{\textcolor[rgb]{0.81,0.36,0.00}{\textbf{#1}}}
\newcommand{\SpecialStringTok}[1]{\textcolor[rgb]{0.31,0.60,0.02}{#1}}
\newcommand{\StringTok}[1]{\textcolor[rgb]{0.31,0.60,0.02}{#1}}
\newcommand{\VariableTok}[1]{\textcolor[rgb]{0.00,0.00,0.00}{#1}}
\newcommand{\VerbatimStringTok}[1]{\textcolor[rgb]{0.31,0.60,0.02}{#1}}
\newcommand{\WarningTok}[1]{\textcolor[rgb]{0.56,0.35,0.01}{\textbf{\textit{#1}}}}

% tightlist command for lists without linebreak
\providecommand{\tightlist}{%
  \setlength{\itemsep}{0pt}\setlength{\parskip}{0pt}}

% From pandoc table feature
\usepackage{longtable,booktabs,array}
\usepackage{calc} % for calculating minipage widths
% Correct order of tables after \paragraph or \subparagraph
\usepackage{etoolbox}
\makeatletter
\patchcmd\longtable{\par}{\if@noskipsec\mbox{}\fi\par}{}{}
\makeatother
% Allow footnotes in longtable head/foot
\IfFileExists{footnotehyper.sty}{\usepackage{footnotehyper}}{\usepackage{footnote}}
\makesavenoteenv{longtable}

% Pandoc citation processing
%From Pandoc 3.1.8
% definitions for citeproc citations
\NewDocumentCommand\citeproctext{}{}
\NewDocumentCommand\citeproc{mm}{%
  \begingroup\def\citeproctext{#2}\cite{#1}\endgroup}
\makeatletter
 % allow citations to break across lines
 \let\@cite@ofmt\@firstofone
 % avoid brackets around text for \cite:
 \def\@biblabel#1{}
 \def\@cite#1#2{{#1\if@tempswa , #2\fi}}
\makeatother
\newlength{\cslhangindent}
\setlength{\cslhangindent}{1.5em}
\newlength{\csllabelwidth}
\setlength{\csllabelwidth}{3em}
\newenvironment{CSLReferences}[2] % #1 hanging-indent, #2 entry-spacing
 {\begin{list}{}{%
  \setlength{\itemindent}{0pt}
  \setlength{\leftmargin}{0pt}
  \setlength{\parsep}{0pt}
  % turn on hanging indent if param 1 is 1
  \ifodd #1
   \setlength{\leftmargin}{\cslhangindent}
   \setlength{\itemindent}{-1\cslhangindent}
  \fi
  % set entry spacing
  \setlength{\itemsep}{#2\baselineskip}}}
 {\end{list}}
\usepackage{calc}
\newcommand{\CSLBlock}[1]{#1\hfill\break}
\newcommand{\CSLLeftMargin}[1]{\parbox[t]{\csllabelwidth}{#1}}
\newcommand{\CSLRightInline}[1]{\parbox[t]{\linewidth - \csllabelwidth}{#1}\break}
\newcommand{\CSLIndent}[1]{\hspace{\cslhangindent}#1}

\begin{document}
\maketitle


\begin{abstract}
How unfortunate is that when you intake drug to cure a specific issue
but end up having other adverse side effects? Adverse Drug Reactions
(ADRs) represent a major challenge in pharmacovigilance and drug safety,
contributing significantly to patient morbidity, mortality, and
healthcare costs. So in this paper, we have evaluated and developed
statistical methods and models in predicting ADRs using two
complementary sources of drug information, drug gene interaction
profiles and chemical fingerprints. The major challenge in working with
ADR dataset is that, there is a huge imbalance in the available dataset,
so the predictive model should understand the imbalance and try to
figure out the pattern.The models evaluated include a Naïve Baseline,
Kernel Regression (KR), Support Vector Machines (SVM) with linear and
RBF kernels, and a novel Weighted Non-negative Matrix Factorisation
(NMF) framework combined with Kernel Regression. Performance is assessed
across 5-fold cross-validation using the Area Under the ROC Curve
(AUROC) and Area Under the Precision-Recall Curve (AUPR). Our results
demonstrate that the Weighted NMF with Kernel Regression achieves the
highest AUPR (0.5007 ± 0.0021), making it the most effective method for
handling the imbalanced nature of ADR datasets, while VKR NMF achieves
the highest AUROC (0.9183 ± 0.0025). We further develop an interactive
web application enabling real time single drug ADR prediction.
\end{abstract}


\section{Introduction}\label{introduction}

Adverse Drug Reactions (ADRs) are unintended, harmful responses to
medications administered at normal therapeutic doses. According to the
WHO, ADRs are among the leading causes of hospitalisation and death
worldwide, accounting for approximately 5--7\% of all hospital
admissions in developed countries. Beyond the human cost, ADRs impose
enormous financial burdens on healthcare systems, with estimates
suggesting costs in the hundreds of billions of dollars annually.

Traditional approaches to ADR detection rely primarily on post market
pharmacovigilance, monitoring drug safety signals after a drug has
already entered clinical use. While these methods are effective, yet
someone has to experience an adverse reaction to flag the behavior.
Hence, computational prediction of ADRs, offers the potential to
anticipate safety issues before clinical exposure, informing drug
development decisions and personalising treatment.

Drugs can be characterised by their interaction profiles with genes and
by their chemical structure, both of which are expected to carry
information about their biological activity and, by extension, their
side effect profiles. However, the construction of such models faces two
fundamental statistical challenges:

\begin{enumerate}
\def\labelenumi{\arabic{enumi}.}
\item
  \textbf{Sparsity}: The drug side effect matrix is extremely sparse.
  Most drug side effect pairs are unobserved, not because the side
  effect is absent, but because it has not been reported or studied.
\item
  \textbf{Class Imbalance}: The number of known positive drug side
  effect associations is vastly smaller than the number of negatives,
  making standard classifiers biased toward predicting the majority
  (negative) class.
\end{enumerate}

Addressing these challenges requires careful model selection and
evaluation. Standard accuracy metrics such as AUROC can be misleading in
highly imbalanced settings. Hence, AUPR is often a more informative
metric when positives are rare.

In this paper, we present a comprehensive study of ADR prediction using
drug gene interaction (DGI) profiles and chemical fingerprints. We
evaluated a range of models from a simple Naïve Baseline to Kernel
Regression, Support Vector Machines, and a novel Weighted Non-negative
Matrix Factorisation (NMF) framework.

\section{Data}\label{data}

\subsection{Dataset}\label{dataset}

The primary data sources were integrated from DGI database,

\begin{itemize}
\item
  \textbf{Drug Gene Interaction (DGI) data}: We retrieved drug gene
  interaction (DGI) pairs from DGI 4.0 database. We mapped the drug
  names from SIDER 4.1 to DGIdb 4.0 and found 973 drugs in common.
\item
  \textbf{Chemical Fingerprints}: We retrieved the drug chemical
  structure fingerprints (Chem) by mapping the CIDs of 1430 drugs from
  SIDER 4.1 to PubChem. PubChem fingerprints are binary vectors that can
  represent the chemical structures of drugs. We used the ChemmineR and
  ChemmineOB packages in Python to generate the fingerprints.
\item
  \textbf{Side Effect data}: We used the SIDER 4.1 database. SIDER 4.1
  is an ADR dataset containing drug ADR data, which provides labels for
  our ADR predictions. We searched features from DGIdb and PubChem for
  these drugs and found 973 common drugs.
\end{itemize}

\subsection{Preprocessing}\label{preprocessing}

All three datasets were binarised (binary matrices) prior to modelling.
The DGI and Chemical fingerprints data had all the drugs in the first
column. Genes and generated fingerprints were presented in each columns
in the respective datasets.

The final dataset comprised:

\begin{itemize}
\tightlist
\item
  \textbf{973 drugs} in total
\item
  \textbf{3,647 unique genes} in the drug-gene interaction matrix
\item
  \textbf{1,024 chemical fingerprint features} per drug
\end{itemize}

Table 1 summarises the dataset dimensions.

\begin{longtable}[]{@{}ll@{}}
\caption{Summary of dataset dimensions.}\tabularnewline
\toprule\noalign{}
Feature Type & Dimensionality \\
\midrule\noalign{}
\endfirsthead
\toprule\noalign{}
Feature Type & Dimensionality \\
\midrule\noalign{}
\endhead
\bottomrule\noalign{}
\endlastfoot
Drug-Gene Interaction & 973 × 3,647 \\
Chemical Fingerprints & 973 × 1,024 \\
Side Effect Matrix & 973 × 5779 \\
\end{longtable}

\subsection{Class Imbalance in the
dataset}\label{class-imbalance-in-the-dataset}

The drug side effect matrix is highly sparse, with a positive rate well
below 10\% for most side effects. This severe class imbalance motivates
the use of AUPR as a primary evaluation metric alongside AUROC, since
AUROC can appear inflated in imbalanced settings. In the entire dataset,
about 97\% of the data represents the negative class and remaining as
positive class

\section{Examples of citations, figures, tables,
references}\label{examples-of-citations-figures-tables-references}

\label{sec:others}

You can insert references. Here is some text (Kour and Saabne 2014b,
2014a) and see Hadash et al. (2018).

The documentation for \verb+natbib+ may be found at

You can use custom blocks with LaTeX support from \textbf{rmarkdown} to
create environment.

\begin{center}
\url{http://mirrors.ctan.org/macros/latex/contrib/natbib/natnotes.pdf\%7D}

\end{center}

Of note is the command \verb+\citet+, which produces citations
appropriate for use in inline text.

You can insert LaTeX environment directly too.

\begin{verbatim}
   \citet{hasselmo} investigated\dots
\end{verbatim}

produces

\begin{quote}
  Hasselmo, et al.\ (1995) investigated\dots
\end{quote}

\begin{center}
  \url{https://www.ctan.org/pkg/booktabs}
\end{center}

\subsection{Figures}\label{figures}

You can insert figure using LaTeX directly.

See Figure \ref{fig:fig1}. Here is how you add footnotes. {[}\^{}Sample
of the first footnote.{]}

\begin{figure}
  \centering
  \fbox{\rule[-.5cm]{4cm}{4cm} \rule[-.5cm]{4cm}{0cm}}
  \caption{Sample figure caption.}
  \label{fig:fig1}
\end{figure}

But you can also do that using R.

You can use \textbf{bookdown} to allow references for Tables and
Figures.

\subsection{Tables}\label{tables}

Below we can see how to use tables.

See awesome Table\textasciitilde{}\ref{tab:table} which is written
directly in LaTeX in source Rmd file.

\begin{table}
 \caption{Sample table title}
  \centering
  \begin{tabular}{lll}
    \toprule
    \multicolumn{2}{c}{Part}                   \\
    \cmidrule(r){1-2}
    Name     & Description     & Size ($\mu$m) \\
    \midrule
    Dendrite & Input terminal  & $\sim$100     \\
    Axon     & Output terminal & $\sim$10      \\
    Soma     & Cell body       & up to $10^6$  \\
    \bottomrule
  \end{tabular}
  \label{tab:table}
\end{table}

You can also use R code for that.

\begin{Shaded}
\begin{Highlighting}[]
\NormalTok{knitr}\SpecialCharTok{::}\FunctionTok{kable}\NormalTok{(}\FunctionTok{head}\NormalTok{(mtcars), }\AttributeTok{caption =} \StringTok{"Head of mtcars table"}\NormalTok{)}
\end{Highlighting}
\end{Shaded}

\begin{longtable}[]{@{}
  >{\raggedright\arraybackslash}p{(\linewidth - 22\tabcolsep) * \real{0.2609}}
  >{\raggedleft\arraybackslash}p{(\linewidth - 22\tabcolsep) * \real{0.0725}}
  >{\raggedleft\arraybackslash}p{(\linewidth - 22\tabcolsep) * \real{0.0580}}
  >{\raggedleft\arraybackslash}p{(\linewidth - 22\tabcolsep) * \real{0.0725}}
  >{\raggedleft\arraybackslash}p{(\linewidth - 22\tabcolsep) * \real{0.0580}}
  >{\raggedleft\arraybackslash}p{(\linewidth - 22\tabcolsep) * \real{0.0725}}
  >{\raggedleft\arraybackslash}p{(\linewidth - 22\tabcolsep) * \real{0.0870}}
  >{\raggedleft\arraybackslash}p{(\linewidth - 22\tabcolsep) * \real{0.0870}}
  >{\raggedleft\arraybackslash}p{(\linewidth - 22\tabcolsep) * \real{0.0435}}
  >{\raggedleft\arraybackslash}p{(\linewidth - 22\tabcolsep) * \real{0.0435}}
  >{\raggedleft\arraybackslash}p{(\linewidth - 22\tabcolsep) * \real{0.0725}}
  >{\raggedleft\arraybackslash}p{(\linewidth - 22\tabcolsep) * \real{0.0725}}@{}}
\caption{Head of mtcars table}\tabularnewline
\toprule\noalign{}
\begin{minipage}[b]{\linewidth}\raggedright
\end{minipage} & \begin{minipage}[b]{\linewidth}\raggedleft
mpg
\end{minipage} & \begin{minipage}[b]{\linewidth}\raggedleft
cyl
\end{minipage} & \begin{minipage}[b]{\linewidth}\raggedleft
disp
\end{minipage} & \begin{minipage}[b]{\linewidth}\raggedleft
hp
\end{minipage} & \begin{minipage}[b]{\linewidth}\raggedleft
drat
\end{minipage} & \begin{minipage}[b]{\linewidth}\raggedleft
wt
\end{minipage} & \begin{minipage}[b]{\linewidth}\raggedleft
qsec
\end{minipage} & \begin{minipage}[b]{\linewidth}\raggedleft
vs
\end{minipage} & \begin{minipage}[b]{\linewidth}\raggedleft
am
\end{minipage} & \begin{minipage}[b]{\linewidth}\raggedleft
gear
\end{minipage} & \begin{minipage}[b]{\linewidth}\raggedleft
carb
\end{minipage} \\
\midrule\noalign{}
\endfirsthead
\toprule\noalign{}
\begin{minipage}[b]{\linewidth}\raggedright
\end{minipage} & \begin{minipage}[b]{\linewidth}\raggedleft
mpg
\end{minipage} & \begin{minipage}[b]{\linewidth}\raggedleft
cyl
\end{minipage} & \begin{minipage}[b]{\linewidth}\raggedleft
disp
\end{minipage} & \begin{minipage}[b]{\linewidth}\raggedleft
hp
\end{minipage} & \begin{minipage}[b]{\linewidth}\raggedleft
drat
\end{minipage} & \begin{minipage}[b]{\linewidth}\raggedleft
wt
\end{minipage} & \begin{minipage}[b]{\linewidth}\raggedleft
qsec
\end{minipage} & \begin{minipage}[b]{\linewidth}\raggedleft
vs
\end{minipage} & \begin{minipage}[b]{\linewidth}\raggedleft
am
\end{minipage} & \begin{minipage}[b]{\linewidth}\raggedleft
gear
\end{minipage} & \begin{minipage}[b]{\linewidth}\raggedleft
carb
\end{minipage} \\
\midrule\noalign{}
\endhead
\bottomrule\noalign{}
\endlastfoot
Mazda RX4 & 21.0 & 6 & 160 & 110 & 3.90 & 2.620 & 16.46 & 0 & 1 & 4 &
4 \\
Mazda RX4 Wag & 21.0 & 6 & 160 & 110 & 3.90 & 2.875 & 17.02 & 0 & 1 & 4
& 4 \\
Datsun 710 & 22.8 & 4 & 108 & 93 & 3.85 & 2.320 & 18.61 & 1 & 1 & 4 &
1 \\
Hornet 4 Drive & 21.4 & 6 & 258 & 110 & 3.08 & 3.215 & 19.44 & 1 & 0 & 3
& 1 \\
Hornet Sportabout & 18.7 & 8 & 360 & 175 & 3.15 & 3.440 & 17.02 & 0 & 0
& 3 & 2 \\
Valiant & 18.1 & 6 & 225 & 105 & 2.76 & 3.460 & 20.22 & 1 & 0 & 3 & 1 \\
\end{longtable}

\subsection{Lists}\label{lists}

\begin{itemize}
\tightlist
\item
  Item 1
\item
  Item 2
\item
  Item 3
\end{itemize}

\protect\phantomsection\label{refs}
\begin{CSLReferences}{1}{1}
\bibitem[\citeproctext]{ref-hadash2018estimate}
Hadash, Guy, Einat Kermany, Boaz Carmeli, Ofer Lavi, George Kour, and
Alon Jacovi. 2018. {``Estimate and Replace: A Novel Approach to
Integrating Deep Neural Networks with Existing Applications.''}
\emph{arXiv Preprint arXiv:1804.09028}.

\bibitem[\citeproctext]{ref-kour2014fast}
Kour, George, and Raid Saabne. 2014a. {``Fast Classification of
Handwritten on-Line Arabic Characters.''} \emph{Soft Computing and
Pattern Recognition (SoCPaR), 2014 6th International Conference of},
312--18.

\bibitem[\citeproctext]{ref-kour2014real}
Kour, George, and Raid Saabne. 2014b. {``Real-Time Segmentation of
on-Line Handwritten Arabic Script.''} \emph{Frontiers in Handwriting
Recognition (ICFHR), 2014 14th International Conference on}, 417--22.

\end{CSLReferences}

\bibliographystyle{unsrt}
\bibliography{references.bib}


\end{document}
